%-------------------------------------------------------------------------------
%	SECTION TITLE
%-------------------------------------------------------------------------------
\cvsection{Extracurricular Activities and Projects}


%-------------------------------------------------------------------------------
%	CONTENT
%-------------------------------------------------------------------------------
\begin{cventries}

%---------------------------------------------------------
  \cventry
    {\underline{\href{github.com/fesoliveira014/voronator-rs}{github.com/fesoliveira014/voronator-rs}}} % Affiliation/role
    {Voronator} % Organization/group
    {} % Location
    {2020 - Present} % Date(s)
    {
      \begin{cvitems}
      \item {Rust crate (library) implementing Voronoi Diagram and Delaunay Triangulation algorithms. It takes in a list of 2D points and returns a graph data structure containing geometrical and spatial information of either just the triangulation or both triangulation and it's dual. It also implements a variant of the Voronoi Diagram that uses centroids instead of circumcenters, allowing for a more aesthetically pleasing result with some trade-offs.}
      \end{cvitems}
    }

  \cventry
    {\underline{\href{https://github.com/fesoliveira014/yanve}{github.com/fesoliveira014/yanve}}} % Affiliation/role
    {Toy Rendering Engine} % Organization/group
    {} % Location
    {2014 - Present} % Date(s)
    {
      \begin{cvitems}
        \item {Since 2014 I have been writing toy rendering and voxel engines for fun. This is the fourth iteration of the engine, providing a wrap over OpenGL, a scene graph and other useful data structures and algorithms for rendering scenes. Past iterations focused almost completely on voxel rendering and experiments on efficient meshing. However, I wanted to make a new version that enabled me to render other things without needing to re-write too much code all the time. It is written in C++17 and using \textit{dear ImGui} for the user interface. The current version does not implement many graphics techniques yet, but past iterations (i.e. \underline{\href{https://github.com/fesoliveira014/playground}{playground}}) implemented techniques such as \textit{Phong Shading} and \textit{Shadow Mapping}. Much inspiration for the design was taken from existing projects, such as magnum and Horde3D.}
      \end{cvitems}
    }

%---------------------------------------------------------
\end{cventries}
